% GENERATED FILE. Edit canonical lessons, then run npm run generate:books.
% canonical-source-hash: fnv1a64:54f3fa434cbd49bc
% canonical-lessons: ML-C489-karansi, ML-C489-vinimayam, ML-C489-cillarayakkuka, ML-C489-panamataykkuka, ML-C489-varinilkkuka

\chapter{Money Matters}
\label{ch:ml-money-matters}
\hlchaptermodality{489}

\begin{chapteropening}
\textbf{By the end of this chapter:} \emph{I can say currency, exchange, to break a note, to make a payment and to stand in line.}

Complete the last lesson of chapter 489: I can say currency, exchange, to break a note, to make a payment and to stand in line.
\end{chapteropening}

\section[\texorpdfstring{\ml{കറൻസി} (kaṟansi)}{കറൻസി (kaṟansi)}]{\ml{കറൻസി} (kaṟansi) — currency}

\label{lesson:ML-C489-karansi}



\begin{quote}
Before the new one: say the Malayalam for a validity period, then the Malayalam for to expire.
\end{quote}

\subsection*{You'll want to know: \ml{കറൻസി}}
\textbf{\ml{കറൻസി}} — \emph{kaṟansi} — \textquotedblleft{}currency\textquotedblright{}.

\begin{quote}
\ml{വിമാനത്താവളത്തിൽ} \ml{വിദേശ} \ml{കറൻസി} \ml{കിട്ടും}.
\end{quote}

\begin{quote}
\emph{vimānattāvaḷattil vidēśa kaṟansi kiṭṭuṁ.}
\end{quote}

\begin{quote}
You can get foreign currency at the airport.
\end{quote}

\ml{വിദേശ} \ml{കറൻസി}, foreign currency; Indian money is the \ml{രൂപ}.

\begin{cousinweb}
\textbf{\ml{കറൻസി}} (\emph{kaṟansi}) was borrowed from English \emph{currency}, the money of a country.
\end{cousinweb}

\begin{grammarlens}[title={What you've built}]
One more word for getting around.
\end{grammarlens}

\subsection*{Your turn}
\begin{itemize}
\raggedright
  \item \emph{Say it:} \emph{kaṟansi}
  \item \emph{Say it:} \emph{kaṟansi}, once more
  \item \emph{Say it:} \emph{vimānattāvaḷattil vidēśa kaṟansi kiṭṭuṁ.} — You can get foreign currency at the airport.
  \item \emph{Recall:} say the Malayalam for a validity period, then the Malayalam for to expire, then say \emph{kaṟansi} again
\end{itemize}

\begin{tcolorbox}[breakable,colback=teal!4,colframe=teal!35!black,title={Before you move on}]
What does \ml{കറൻസി} mean? (\textquotedblleft{}Currency\textquotedblright{}.) Say it once more.
\end{tcolorbox}
\section[\texorpdfstring{\ml{വിനിമയം} (vinimayaṁ)}{വിനിമയം (vinimayaṁ)}]{\ml{വിനിമയം} (vinimayaṁ) — exchange}

\label{lesson:ML-C489-vinimayam}



\begin{quote}
Before the new one: say the Malayalam for to expire, then the Malayalam for currency.
\end{quote}

\subsection*{You'll want to know: \ml{വിനിമയം}}
\textbf{\ml{വിനിമയം}} — \emph{vinimayaṁ} — \textquotedblleft{}exchange\textquotedblright{}.

\begin{quote}
\ml{കറൻസി} \ml{വിനിമയം} \ml{ബാങ്കിൽ} \ml{ചെയ്യാം}.
\end{quote}

\begin{quote}
\emph{kaṟansi vinimayaṁ bāṅkil ceyyāṁ.}
\end{quote}

\begin{quote}
Currency exchange can be done at the bank.
\end{quote}

The exchange rate is the \ml{വിനിമയ} \ml{നിരക്ക്}.

\begin{cousinweb}
\textbf{\ml{വിനിമയം}} (\emph{vinimayaṁ}) is borrowed from Sanskrit \emph{vinimaya}, \textquotedblleft{}exchange, barter\textquotedblright{} (Monier-Williams).
\end{cousinweb}

\begin{grammarlens}[title={What you've built}]
One more word for getting around.
\end{grammarlens}

\subsection*{Your turn}
\begin{itemize}
\raggedright
  \item \emph{Say it:} \emph{vinimayaṁ}
  \item \emph{Say it:} \emph{vinimayaṁ}, once more
  \item \emph{Say it:} \emph{kaṟansi vinimayaṁ bāṅkil ceyyāṁ.} — Currency exchange can be done at the bank.
  \item \emph{Recall:} say the Malayalam for to expire, then the Malayalam for currency, then say \emph{vinimayaṁ} again
\end{itemize}

\begin{tcolorbox}[breakable,colback=teal!4,colframe=teal!35!black,title={Before you move on}]
What does \ml{വിനിമയം} mean? (\textquotedblleft{}Exchange\textquotedblright{}.) Say it once more.
\end{tcolorbox}
\section[\texorpdfstring{\ml{ചില്ലറയാക്കുക} (cillaṟayākkuka)}{ചില്ലറയാക്കുക (cillaṟayākkuka)}]{\ml{ചില്ലറയാക്കുക} (cillaṟayākkuka) — to break a note}

\label{lesson:ML-C489-cillarayakkuka}



\begin{quote}
Before the new one: say the Malayalam for currency, then the Malayalam for exchange.
\end{quote}

\subsection*{You'll want to know: \ml{ചില്ലറയാക്കുക}}
\textbf{\ml{ചില്ലറയാക്കുക}} — \emph{cillaṟayākkuka} — \textquotedblleft{}to break a note\textquotedblright{}.

\begin{quote}
\ml{കടക്കാരൻ} \ml{എന്റെ} \ml{അഞ്ഞൂറിന്റെ} \ml{നോട്ട്} \ml{ചില്ലറയാക്കി}.
\end{quote}

\begin{quote}
\emph{kaṭakkāran enṟe aññūṟinṟe nōṭṭŭ cillaṟayākki.}
\end{quote}

\begin{quote}
The shopkeeper broke my five-hundred note.
\end{quote}

Past: \ml{ചില്ലറയാക്കി}, \emph{cillaṟayākki}.

Literally to make into small change: \ml{ചില്ലറ} plus \ml{ആക്കുക}, to make.

\begin{cousinweb}
\textbf{\ml{ചില്ലറയാക്കുക}} (\emph{cillaṟayākkuka}) is a Malayalam verb built from the native word \ml{ചില്ലറ}, small change (DEDR 2574), and \ml{ആക്കുക}, to make.
\end{cousinweb}

\begin{grammarlens}[title={What you've built}]
One more word for getting around.
\end{grammarlens}

\subsection*{Your turn}
\begin{itemize}
\raggedright
  \item \emph{Say it:} \emph{cillaṟayākkuka}
  \item \emph{Say it:} \emph{cillaṟayākkuka}, once more
  \item \emph{Say it:} \emph{kaṭakkāran enṟe aññūṟinṟe nōṭṭŭ cillaṟayākki.} — The shopkeeper broke my five-hundred note.
  \item \emph{Recall:} say the Malayalam for currency, then the Malayalam for exchange, then say \emph{cillaṟayākkuka} again
\end{itemize}

\begin{tcolorbox}[breakable,colback=teal!4,colframe=teal!35!black,title={Before you move on}]
What does \ml{ചില്ലറയാക്കുക} mean? (\textquotedblleft{}To break a note\textquotedblright{}.) Say it once more.
\end{tcolorbox}
\section[\texorpdfstring{\ml{പണമടയ്ക്കുക} (paṇamaṭaykkuka)}{പണമടയ്ക്കുക (paṇamaṭaykkuka)}]{\ml{പണമടയ്ക്കുക} (paṇamaṭaykkuka) — to make a payment; to pay}

\label{lesson:ML-C489-panamataykkuka}



\begin{quote}
Before the new one: say the Malayalam for exchange, then the Malayalam for to break a note.
\end{quote}

\subsection*{You'll want to know: \ml{പണമടയ്ക്കുക}}
\textbf{\ml{പണമടയ്ക്കുക}} — \emph{paṇamaṭaykkuka} — \textquotedblleft{}to make a payment; to pay\textquotedblright{}.

\begin{quote}
\ml{മുറി} \ml{ഒഴിയുന്നതിന്} \ml{മുൻപ്} \ml{ഞങ്ങൾ} \ml{പണമടച്ചു}.
\end{quote}

\begin{quote}
\emph{muṟi oḻiyunnatinŭ munpŭ ñaṅṅaḷ paṇamaṭaccu.}
\end{quote}

\begin{quote}
We paid before we vacated the room.
\end{quote}

Past: \ml{പണമടച്ചു}, \emph{paṇamaṭaccu}.

\ml{പണം}, money, plus \ml{അടയ്ക്കുക}, which means both to shut and to pay in: a bill, a fee, a tax.

\begin{cousinweb}
\textbf{\ml{പണമടയ്ക്കുക}} (\emph{paṇamaṭaykkuka}) is a Malayalam compound built from \ml{പണം}, money, from Sanskrit \emph{paṇa}, a coin (CDIAL 7715), and the native verb \ml{അടയ്ക്കുക}, to shut (DEDR 83), which in Malayalam also means to pay in.
\end{cousinweb}

\begin{grammarlens}[title={What you've built}]
One more word for getting around.
\end{grammarlens}

\subsection*{Your turn}
\begin{itemize}
\raggedright
  \item \emph{Say it:} \emph{paṇamaṭaykkuka}
  \item \emph{Say it:} \emph{paṇamaṭaykkuka}, once more
  \item \emph{Say it:} \emph{muṟi oḻiyunnatinŭ munpŭ ñaṅṅaḷ paṇamaṭaccu.} — We paid before we vacated the room.
  \item \emph{Recall:} say the Malayalam for exchange, then the Malayalam for to break a note, then say \emph{paṇamaṭaykkuka} again
\end{itemize}

\begin{tcolorbox}[breakable,colback=teal!4,colframe=teal!35!black,title={Before you move on}]
What does \ml{പണമടയ്ക്കുക} mean? (\textquotedblleft{}To make a payment; to pay\textquotedblright{}.) Say it once more.
\end{tcolorbox}
\section[\texorpdfstring{\ml{വരിനിൽക്കുക} (varinilkkuka)}{വരിനിൽക്കുക (varinilkkuka)}]{\ml{വരിനിൽക്കുക} (varinilkkuka) — to stand in line; to queue}

\label{lesson:ML-C489-varinilkkuka}



\begin{quote}
Before the new one: say the Malayalam for to break a note, then the Malayalam for to make a payment.
\end{quote}

\subsection*{You'll want to know: \ml{വരിനിൽക്കുക}}
\textbf{\ml{വരിനിൽക്കുക}} — \emph{varinilkkuka} — \textquotedblleft{}to stand in line; to queue\textquotedblright{}.

\begin{quote}
\ml{ടിക്കറ്റിനായി} \ml{ഞങ്ങൾ} \ml{ഒരു} \ml{മണിക്കൂർ} \ml{വരിനിന്നു}.
\end{quote}

\begin{quote}
\emph{ṭikkaṟṟināyi ñaṅṅaḷ oru maṇikkūr varininnu.}
\end{quote}

\begin{quote}
We queued for an hour for tickets.
\end{quote}

Past: \ml{വരിനിന്നു}, \emph{varininnu}.

\ml{വരി}, a line or row, plus \ml{നിൽക്കുക}, to stand. A queue is also a \ml{ക്യൂ}.

\begin{cousinweb}
\textbf{\ml{വരിനിൽക്കുക}} (\emph{varinilkkuka}) is a native Dravidian compound of \ml{വരി}, a row or line (DEDR 5263), and \ml{നിൽക്കുക}, to stand (DEDR 3675).
\end{cousinweb}

\begin{grammarlens}[title={What you've built}]
One more word for getting around.
\end{grammarlens}

\subsection*{Your turn}
\begin{itemize}
\raggedright
  \item \emph{Say it:} \emph{varinilkkuka}
  \item \emph{Say it:} \emph{varinilkkuka}, once more
  \item \emph{Say it:} \emph{ṭikkaṟṟināyi ñaṅṅaḷ oru maṇikkūr varininnu.} — We queued for an hour for tickets.
  \item \emph{Recall:} say the Malayalam for to break a note, then the Malayalam for to make a payment, then say \emph{varinilkkuka} again
\end{itemize}

\begin{tcolorbox}[breakable,colback=teal!4,colframe=teal!35!black,title={Before you move on}]
What does \ml{വരിനിൽക്കുക} mean? (\textquotedblleft{}To stand in line; to queue\textquotedblright{}.) Say it once more.
\end{tcolorbox}
